\chapter{АНАЛИЗ СОСТОЯНИЯ ВОПРОСА И ПОСТАНОВКА ЗАДАЧ ПРОЕКТИРОВАНИЯ}

\section{Обзор существующих методов и технических решений}

На современном этапе развития технологий автоматизации мониторинг параметров окружающей среды переходит от проводных интерфейсов к беспроводным решениям на базе стека протоколов TCP/IP. Согласно литературным источникам \cite{ivanov2023}, наиболее критичными параметрами для серверных помещений являются температура ($T$) и относительная влажность ($H$).

Для измерения данных параметров традиционно применяются следующие типы датчиков:
\begin{enumerate}
    \item полупроводниковые (например, серия DHT);
    \item емкостные (SHT3x);
    \item резистивные.
\end{enumerate}

Сравнительная характеристика популярных датчиков приведена в таблице 1.1.

\begin{table}[h!]
    \centering
    \caption{Сравнительные характеристики первичных преобразователей}
    \begin{tabular}{|l|c|c|c|}
        \hline
        Модель датчика & Диапазон, $^\circ$C & Погрешность, \% & Интерфейс \\ \hline
        DHT11          & 0...50              & $\pm$ 5        & Single-wire \\ \hline
        DHT22          & -40...80            & $\pm$ 2        & Single-wire \\ \hline
        SHT31          & -40...125           & $\pm$ 0.3      & I2C \\ \hline
    \end{tabular}
\end{table}

\section{Анализ сетевых протоколов передачи данных}

Для передачи телеметрической информации в сетях IoT наиболее эффективным признан протокол MQTT (Message Queuing Telemetry Transport). Его основным преимуществом перед HTTP является низкий оверхед (минимальный размер заголовка составляет 2 байта), что критично для устройств с батарейным питанием.

\section{Обоснование выбора направления проектирования}

На основании проведенного анализа источников литературы и патентного поиска можно сделать вывод, что для реализации поставленной цели оптимальным является использование микроконтроллера с интегрированным модулем Wi-Fi и датчика цифрового типа с интерфейсом I2C. 

Это позволит обеспечить:
\begin{itemize}
    \item высокую точность измерений;
    \item возможность удаленного обновления прошивки (OTA);
    \item масштабируемость системы за счет использования брокера сообщений.
\end{itemize}

\newpage